% Safety-Aware Speculative Decoding — NeurIPS-style Project Report
% Page limit: 9 pages including references
% Format: thick rules above/below title block, centered title (sans-serif), centered author.

\documentclass[11pt,a4paper]{article}
\usepackage[utf8]{inputenc}
\usepackage[T1]{fontenc}
\usepackage{hyperref}
\usepackage{url}
\usepackage{booktabs}
\usepackage{amsfonts,amsmath,amssymb}
\usepackage{graphicx}
\usepackage{natbib}
\usepackage{algorithm}
\usepackage{algorithmic}
\usepackage{xcolor}

% If using neurips_2024.sty in Overleaf, uncomment as needed:
% \usepackage{neurips_2024}

\begin{document}

% --- Thick rule above title ---
\noindent\rule{\textwidth}{1.5pt}\\[1em]
% --- Title: centered, large, bold, sans-serif ---
\begin{center}
  {\sffamily\bfseries\LARGE Safety-Aware Speculative Decoding:\\
  Extending SSD with Composite Risk Scores and Hidden-State Steering}\\[1.5em]
  % --- Author block: centered, serif ---
  Raghusrinivasan Venkatesan\\
  UC San Diego\\
  \texttt{rlvenkatesan@ucsd.edu}\\
  % Uncomment if needed: 2 Grace days used
\end{center}
% --- Thick rule below author ---
\vspace{0.5em}
\noindent\rule{\textwidth}{1.5pt}
\vspace{1em}

% --- Abstract ---
\begin{abstract}
\noindent
Large language models (LLMs) remain vulnerable to jailbreak attacks despite alignment. We build on \emph{Speculative Safety-Aware Decoding} (SSD), which uses a small safety-conscious draft model and match-ratio-based mode switching to improve safety and inference speed. We extend SSD along three directions: (1) a \emph{Composite Risk Score} (CRS) that replaces coarse binary switching with a per-token continuous risk signal combining match mismatch, KL divergence, entropy gap, and refusal-token mass; (2) \emph{hidden-state steering} that injects the target model's refusal-aligned representations into the draft, with injection scale modulated by CRS; and (3) integration with \emph{Alignment-Augmented Speculative Decoding} (AASD) for KV-cached generation and alignment sampling. We also add training-free extensions: contrastive system prompts and perplexity-gated early Union mode. On a custom high-ASR evaluation set (70 harmful, 30 benign prompts), CRS alone reduces attack success rate (ASR) by 30.7 percentage points versus vanilla; CRS combined with steering achieves \textbf{1.8\%} overall ASR and halves DeepInception ASR (40\% $\to$ 20\%) with minimal over-refusal (3.3\%). We ablate CRS vs.\ steering and discuss failure modes on template-based jailbreaks.
\end{abstract}

\section{Introduction}
\label{sec:intro}

Aligned large language models (LLMs) are routinely deployed for assistive tasks, yet they remain susceptible to \emph{jailbreak} attacks that elicit harmful or unsafe behavior through adversarial prompts, roleplay framing, or nested narrative techniques~\cite{deepinception,jbb}. Defenses that rely on retraining or heavy fine-tuning are costly and may not generalize to new attack types. \emph{Inference-time} interventions—operating only during decoding—offer a lightweight alternative: they require no changes to the target model's weights and can be combined with speculative decoding to simultaneously improve safety and speed.

\emph{Speculative decoding} uses a small draft model to propose tokens that a larger target model verifies in parallel, reducing latency~\cite{leviathan2023,specdecode}. \citet{ssd2025} introduced \textbf{Speculative Safety-Aware Decoding (SSD)}, which repurposes this setup for safety: a safety-conscious draft model and the target model produce token distributions at each step, and a \emph{match ratio} over a fixed window (e.g., 7 tokens) determines whether to sample from the \emph{intersection} or \emph{union} of their top-$c$ token sets. High match ratio is interpreted as benign; low match triggers a more permissive \emph{Union} mode that allows the draft's refusal tokens into the sample space. SSD improves safety while retaining inference speedup, but its binary, window-based switching is coarse and can miss early-token jailbreak signals or over-refuse on long benign responses.

In this work we extend SSD along several directions implemented and evaluated in a single project pipeline:

\begin{itemize}
  \item \textbf{Composite Risk Score (CRS):} We replace the binary match-ratio switch with a continuous \emph{Composite Risk Score} computed every token, combining (i) match mismatch, (ii) KL divergence between draft and target distributions, (iii) entropy gap, and (iv) refusal-token mass. Mode switching and steering strength are driven by CRS, enabling finer-grained safety control.
  \item \textbf{Hidden-state steering:} We inject the target model's last hidden state (via a learned linear projection) into the draft model's MLP layers, scaling the injection by CRS so that steering amplifies when risk is high. The projection is trained with \emph{refusal distillation} on harmful prompts where the target refuses.
  \item \textbf{Alignment-Augmented Speculative Decoding (AASD):} We integrate SSD with AASD-style \emph{alignment sampling} (draft logits blended with target prefill prior) and KV-cache-based incremental decoding for O(1) per-step cost, improving both safety alignment and throughput.
  \item \textbf{Training-free extensions:} We use \emph{contrastive system prompts} (safety-heavy for draft, standard for target) and \emph{perplexity-gated early Union}: if prompt perplexity under the draft exceeds a threshold, Union mode is forced from token 1 to handle GCG-style and other high-ppl attacks.
\end{itemize}

We evaluate on a custom dataset of 70 harmful prompts (40 DeepInception, 30 JailbreakBench-wrapped) and 30 benign XSTest prompts, with Qwen3Guard as the safety classifier. Our main findings: (1) CRS is the primary driver of safety gain (30.7 pp ASR reduction vs.\ vanilla); (2) steering alone with vanilla SSD mode switching gives limited benefit; (3) CRS + steering achieves 1.8\% overall ASR and halves DeepInception ASR versus CRS-only, with 3.3\% over-refusal; (4) template-wrapped jailbreaks can slightly increase ASR under SSD variants—an important failure mode.

The rest of the paper is organized as follows. \S\ref{sec:related} briefly reviews related work. \S\ref{sec:prelim} formalizes the problem and SSD baseline. \S\ref{sec:method} describes CRS, steering, AASD, and training-free extensions. \S\ref{sec:experiments} and \S\ref{sec:results} present setup and results. \S\ref{sec:discussion} discusses limitations and \S\ref{sec:conclusion} concludes.

\section{Related Work}
\label{sec:related}

\textbf{Jailbreak defenses.} Defenses range from input filtering~\cite{smoothllm}, paraphrase or repetition checks~\cite{parden}, to inference-time decoding strategies. SafeDecoding~\cite{safedecoding} adjusts the sampling distribution using a safety expert model (e.g., CensorTune) by blending base and expert logits and sampling from the intersection of top-$k$ tokens; it does not use speculative decoding. Our baseline SSD~\cite{ssd2025} instead uses a single draft model for both speed and safety and switches between intersection/union based on match ratio.

\textbf{Speculative decoding.} Leviathan et al.~\cite{leviathan2023} and Chen et al.~\cite{specdecode} introduced speculative decoding with a draft model proposing multiple tokens and the target verifying in parallel. Extensions include adaptive draft length~\cite{adaedl,adasd}, entropy-aware rejection~\cite{easd}, and risk-bounded acceptance~\cite{arcdecode}. AASD-style alignment sampling blends draft logits with the target's prefill prior and uses KV caches for O(1) per-step decoding; we integrate this with SSD.

\textbf{Representation steering.} Activation addition~\cite{actadd} steers models at inference by adding contrastive activation differences. Steering pretrained drafters during speculative decoding~\cite{steeringdraft} injects target information into the draft; we adopt a similar idea with a CRS-modulated refusal-distillation projection.

\section{Problem Setup and SSD Baseline}
\label{sec:prelim}

\subsection{Setting}
We have a \emph{target} LLM $P_{\mathrm{t}}$ (e.g., Qwen2.5-7B-Instruct) and a smaller \emph{draft} LLM $P_{\mathrm{d}}$ (e.g., Qwen2.5-1.5B-Instruct). Both are given the same user prompt (possibly with different system prompts). At each decoding step $t$, they produce token distributions $p_{\mathrm{t}}(\cdot\mid x_{<t})$ and $p_{\mathrm{d}}(\cdot\mid x_{<t})$. The goal is to produce a sequence that (i) stays close to the target's behavior on benign prompts (utility) and (ii) refuses or avoids harmful content on jailbreak prompts (safety), without retraining the target.

\subsection{SSD Algorithm (Baseline)}
\label{sec:ssd}

SSD~\cite{ssd2025} uses the draft as a safety-conscious partner. At each step:
\begin{enumerate}
  \item Compute top-$c$ token sets $\mathcal{T}_{\mathrm{t}}$, $\mathcal{T}_{\mathrm{d}}$ from the two distributions.
  \item Maintain a \emph{match ratio} $\beta$ over the last $b$ tokens: the fraction of steps where the \emph{sampled} token was in the draft's top-$c$.
  \item If $\beta > \beta_{\mathrm{th}}$: \emph{Intersection} mode — sample from $\mathcal{T}_{\mathrm{t}} \cap \mathcal{T}_{\mathrm{d}}$ with composite score $p_{\mathrm{t}} + \alpha_I (p_{\mathrm{d}} - p_{\mathrm{t}})$.
  \item If $\beta \le \beta_{\mathrm{th}}$: \emph{Union} mode — sample from $\mathcal{T}_{\mathrm{t}} \cup \mathcal{T}_{\mathrm{d}}$ with $p_{\mathrm{t}} + \alpha_U (p_{\mathrm{d}} - p_{\mathrm{t}})$, allowing the draft's refusal tokens into the support.
\end{enumerate}
Annealing (Algorithm 2 in the original paper) decays $\beta_{\mathrm{th}}$ and $\alpha_I$ when the scheme stays unchanged to limit over-refusal on long benign replies. Our implementation uses contrastive system prompts (draft: safety-heavy; target: standard) and an optional \emph{perplexity gate}: if the prompt's perplexity under the draft exceeds a threshold, we force Union from step 1 to handle adversarial suffixes.

\subsection{Limitations of Binary Match-Ratio Switching}
The match ratio is computed every $b$ tokens and is binary in effect (above or below threshold). Thus (i) the first $b$ tokens have no safety signal from agreement, (ii) slow-building jailbreaks (e.g., narrative) may keep $\beta$ high until harm is already underway, and (iii) the signal is an average over the window rather than a per-token risk. Our CRS addresses these by providing a continuous, per-token risk that can drive both mode choice and steering strength.

\section{Method}
\label{sec:method}

\subsection{Composite Risk Score (CRS)}
\label{sec:crs}

At each token step $t$, we define a scalar \textbf{Composite Risk Score} $r_t \in [0,1]$ from four normalized signals:

\begin{itemize}
  \item \textbf{Match mismatch} $1 - \mathbb{I}[\hat{x}_t \in \mathcal{T}_{\mathrm{d}}]$: disagreement between the (previously) sampled token and the draft's top-$c$ (we use a running indicator).
  \item \textbf{KL divergence} $\mathrm{KL}(p_{\mathrm{d}} \| p_{\mathrm{t}})$: large values indicate the draft distribution diverges from the target (e.g., draft refuses, target does not).
  \item \textbf{Entropy gap} $H(p_{\mathrm{t}}) - H(p_{\mathrm{d}})$: draft more confident than target can indicate an unusual (e.g., harmful) context.
  \item \textbf{Refusal mass} $\sum_{v \in \mathcal{R}} p_{\mathrm{d}}(v)$: probability the draft assigns to a fixed set of refusal tokens (e.g., ``I cannot'', ``I'm sorry'').
\end{itemize}

We normalize each component to $[0,1]$ (e.g., min-max over a short history or clip) and take a convex combination:
\begin{equation}
  r_t = w_1 (1-\mathrm{match}) + w_2 \,\overline{\mathrm{KL}} + w_3 \,\overline{\Delta H} + w_4 \,\overline{\mathrm{refusal}},
\end{equation}
with $w_1{=}w_2{=}0.3$, $w_3{=}0.1$, $w_4{=}0.3$ in our experiments. The score is smoothed over a window of 3 steps. When $\bar{r}_t > \tau_{\mathrm{crs}}$ (e.g., 0.2), we apply stronger safety (intersection with annealing); otherwise we use union. Thus CRS replaces the 7-token binary switch with a per-token, continuous gate.

\subsection{Hidden-State Steering}
\label{sec:steering}

We add a \emph{steering} module that injects the target's last-layer hidden state into the draft's MLP layers so the draft's internal representations move toward the target's refusal behavior when risk is high.

\textbf{Projection.} Let $h_{\mathrm{t}} \in \mathbb{R}^{d_{\mathrm{t}}}$ be the target's last hidden state and $d_{\mathrm{d}}$ the draft's hidden size. We learn a linear map $\mathrm{proj}: \mathbb{R}^{d_{\mathrm{t}}} \to \mathbb{R}^{d_{\mathrm{d}}}$ by \emph{refusal distillation}: on harmful prompts where the target refuses, we minimize the KL between the draft's next-token distribution (after receiving a projected vector) and the target's, or a hidden-state loss (e.g., MSE between draft hidden states and projected target hidden states). Training uses only harmful/refusal data; training on general helpfulness data was found to increase ASR.

\textbf{Inference.} We unit-normalize $h_{\mathrm{t}}$ (critical: unnormalized norms $\approx 295$ cause logit collapse), compute $g_t = \mathrm{proj}(h_{\mathrm{t}}/\|h_{\mathrm{t}}\|)$, and add $g_t$ scaled by $\eta_t$ to the draft's MLP outputs at all layers (or a fraction of layers). The scale is CRS-modulated:
\begin{equation}
  \eta_t = \eta_0 \bigl(1 + \lambda \,\bar{r}_t\bigr),
\end{equation}
so steering is stronger when CRS is high. We use $\eta_0 = 0.003$, $\lambda = 5$ in experiments.

\subsection{Integration with AASD}
\label{sec:aasd}

We combine SSD with an AASD-style pipeline for efficiency and alignment:

\begin{enumerate}
  \item \textbf{Prefill \& KV cache:} Run a single full forward pass for both models on the prompt; keep target and draft KV caches and the target's \emph{prefill logits} at the last prompt position.
  \item \textbf{Alignment sampling:} At each step, draft logits are blended with the (fixed) target prefill prior: $\ell_{\mathrm{d}} \leftarrow (1-\lambda_{\mathrm{align}}) \ell_{\mathrm{d}} + \lambda_{\mathrm{align}} \,\ell_{\mathrm{t}}^{\mathrm{prefill}}$, biasing the draft toward the target's initial response direction.
  \item \textbf{Incremental decoding:} For each new token, only the new token is fed in; KV caches are updated (O(1) per step instead of O(n) full recompute).
  \item \textbf{Conditional verification:} In intersection mode, if $p_{\mathrm{t}}(\hat{x}_t) < \tau \cdot \max_v p_{\mathrm{t}}(v)$, we discard the token and resample from the union set with $\alpha_U$, avoiding low-probability target tokens that could be harmful.
\end{enumerate}

SSD's intersection/union and match-ratio (or CRS) logic run on top of this pipeline; the draft logits used for token sets and composite sampling are the aligned draft logits above.

\subsection{Training-Free Extensions}
\label{sec:trainingfree}

\begin{itemize}
  \item \textbf{Contrastive system prompts:} The draft receives a strong safety-focused system prompt; the target receives a standard ``helpful assistant'' prompt. This widens the draft-target distribution gap on harmful prompts without fine-tuning, making match ratio (and CRS) more discriminative.
  \item \textbf{Perplexity-gated early Union:} Before generation, we compute prompt perplexity under the draft. If it exceeds a threshold (e.g., 50 or 150), we force Union mode from token 1 to handle GCG-style and other high-perplexity jailbreaks where the match-ratio window would otherwise be inactive for the first $b$ tokens.
\end{itemize}

\section{Experimental Setup}
\label{sec:experiments}

\textbf{Models.} Target: Qwen/Qwen2.5-7B-Instruct; draft: Qwen/Qwen2.5-1.5B-Instruct; guard: Qwen/Qwen3Guard-Gen-0.6B. All models 4-bit NF4 quantized. Max new tokens: 256. Hardware: 2$\times$ NVIDIA RTX 4000 Ada (20GB each).

\textbf{Datasets.} We use a custom evaluation set to obtain meaningful ASR spread: 70 harmful prompts (40 DeepInception-style nested fiction + 30 JailbreakBench behaviors wrapped in DAN/AIM/roleplay-style templates) and 30 benign prompts from XSTest (surface-dangerous wording, benign intent) to measure over-refusal. All safety labels are from Qwen3Guard (no keyword heuristics).

\textbf{Methods.} (1) \textbf{Vanilla}: target only, greedy decoding. (2) \textbf{SSD}: match-ratio switching every $b{=}7$, contrastive prompts, PPL gate. (3) \textbf{SSD-CRS}: CRS-driven mode switching, no steering. (4) \textbf{SSD-Steering}: CRS + hidden-state steering (projection trained with refusal distillation). (5) \textbf{SSD-Steering-Only}: same steering and projection as SSD-Steering but vanilla SSD (binary match ratio) and fixed steering scale—ablation to isolate CRS vs.\ steering.

\textbf{Metrics.} Attack Success Rate (ASR) on harmful prompts (lower is better); over-refusal rate on benign XSTest (lower is better). We also report per-source ASR (DeepInception, JBB-wrapped) where available.

\section{Results}
\label{sec:results}

Table~\ref{tab:main} summarizes the main results. Vanilla ASR is 34.3\% overall and 57.5\% on DeepInception. SSD (binary match ratio) reduces overall ASR to 28.6\% ($-5.7$ pp). Replacing the binary switch with CRS (SSD-CRS) yields a large drop to 3.6\% ($-30.7$ pp vs.\ vanilla). Adding hidden-state steering on top of CRS (SSD-Steering) further reduces overall ASR to \textbf{1.8\%} and DeepInception ASR from 40\% to \textbf{20\%}. Over-refusal is 0\% for vanilla and SSD, and 3.3\% (1/30 benign) for CRS and SSD-Steering.

\begin{table}[t]
  \centering
  \caption{Attack Success Rate (ASR) and over-refusal. Lower ASR is better. ``—'' indicates not evaluated on that split.}
  \label{tab:main}
  \begin{tabular}{lcccc}
    \toprule
    Method & Overall ASR$\downarrow$ & DeepInception ASR & JBB Wrapped ASR & Over-refusal$\downarrow$ \\
    \midrule
    Vanilla & 34.3 & 57.5 & 3.3 & 0.0 \\
    SSD & 28.6 & 42.5 & 10.0 & 0.0 \\
    SSD-Steering-Only & 30.0 & 42.5 & 13.3 & 0.0 \\
    SSD-CRS & 3.6 & 40.0 & — & 3.3 \\
    \textbf{SSD-Steering} & \textbf{1.8} & \textbf{20.0} & — & 3.3 \\
    \bottomrule
  \end{tabular}
\end{table}

\textbf{Ablation: CRS vs.\ steering.} SSD-Steering-Only (steering + vanilla match-ratio, no CRS) has 30.0\% ASR—only slightly better than SSD (28.6\%) and much worse than SSD-CRS (3.6\%). So \emph{CRS is the primary driver} of the safety gain; steering alone does not get a strong signal from binary 7-token switching. When CRS is used, adding steering gives an extra $\sim$2 pp overall and halves DeepInception ASR (40\% $\to$ 20\%), indicating that steering helps most when triggered by a fine-grained risk signal.

\textbf{Per-source.} DeepInception is the hardest category (57.5\% vanilla). CRS alone brings it to 40\%; CRS+steering to 20\%. On JBB-wrapped prompts, vanilla ASR is already low (3.3\%); SSD variants slightly \emph{increase} ASR (10--13\%), suggesting template wrappers can confuse the draft's safety signal and lead to overly permissive mode choices—a failure mode for future work.

\section{Discussion}
\label{sec:discussion}

\textbf{Why CRS beats match ratio.} The 7-token match ratio is a delayed, averaged signal. CRS reacts every token and combines multiple views (disagreement, KL, entropy, refusal mass), so it can trigger safety mode earlier and with fewer false negatives on narrative jailbreaks.

\textbf{Why steering needs CRS.} With binary switching, the draft is either steered at a fixed scale or not; the switch happens only every $b$ steps. CRS provides a continuous, per-token gain for the steering strength, so the draft is pushed toward refusal precisely when the risk score is high.

\textbf{Limitations.} (1) Steering uses a trained projection (refusal distillation); the rest of the pipeline is training-free except for that component. (2) JBB-style wrappers can hurt SSD variants; combining with input-level or embedding-level detectors may help. (3) Our dataset is a single custom split; broader benchmarks would strengthen claims.

\section{Conclusion}
\label{sec:conclusion}

We extended Speculative Safety-Aware Decoding (SSD) with a per-token Composite Risk Score (CRS), CRS-modulated hidden-state steering, and integration with AASD for efficient KV-cached decoding. On a high-ASR evaluation set, CRS alone reduces attack success rate by 30.7 pp; CRS plus steering achieves 1.8\% overall ASR and halves DeepInception ASR with minimal over-refusal. Ablations show that CRS is necessary for steering to be effective. We also identified a failure mode on template-wrapped jailbreaks. Future work may combine CRS with Jensen-Shannon-based acceptance signals and adaptive draft length for joint safety and speedup gains.

\small
\bibliographystyle{plainnat}
\bibliography{refs}
\end{document}
